\subsectionfit{Crontab}
\index{crontab}
\index{Config!crontab}

In configuration, some fields define a crontab which defines when something should happen over time.
This is a simple string which consists of either a pre-defined name or a definition.

\subsubsection{Crontab definitions}

Crontab definitions follow one of two formats with fields defined in \vreftab{config:crontab:fields}:

\textbf{``minute hour dom month dow''} which follows the standard Unix cron format.

\textbf{``second minute hour dom month dow''} with an additional field representing seconds at the start.

\begin{table}[ht]
 \caption{Crontab field definitions}
 \labtab{config:crontab:fields}
 \begin{tabular}{l l l}
 \toprule
 field & allowed values \\
 \midrule
 second & 0\dots59 & Optional, not used in unix cron\\
 minute & 0\dots59 \\
 hour & 0\dots23 \\
 dom & 0\dots31 & day of month \\
 month & 0\dots12 \\
 dow & 0\dots7 & day of week, 0 or 7 is Sunday\\
 \bottomrule
 \end{tabular}
\end{table}

A field may be an asterisk (*), which always stands for ``first-last''.

Ranges of numbers are allowed.
Ranges are two numbers separated with a hyphen.
The specified range is inclusive.
For example, 8--11 for an ``hours'' entry specifies execution at hours 8, 9, 10 and 11.

Lists are allowed.
A list is a set of numbers (or ranges) separated by commas.
Examples: ``1,2,5,9'', ``0-4,8-12''.

Step values can be used in conjunction with ranges.
Following a range with ``/<number>'' specifies skips of the number's value through the range.
For example, ``0-23/2'' can be used in the hours field to specify command execution every other hour.
Steps are also permitted after an asterisk, so if you want to say ``every two hours'', just use ``*/2''.

Names can also be used for the ``month'' and ``day of week'' fields.
Use the first three letters of the particular day or month (case doesn't matter).
Ranges or lists of names are not allowed.
